\documentclass[12pt]{article}
\usepackage{amsmath,amssymb,amsthm}
\usepackage{geometry}
\geometry{margin=1in}
\usepackage{hyperref}
\usepackage{booktabs}
\usepackage{authblk}
\usepackage[T1]{fontenc}

\newtheorem{theorem}{Theorem}
\newtheorem{proposition}[theorem]{Proposition}
\newtheorem{corollary}[theorem]{Corollary}
\theoremstyle{definition}
\newtheorem{definition}[theorem]{Definition}
\newtheorem{remark}[theorem]{Remark}

\title{Neutrino Masses from the Golden Ratio Lattice:\\
A Parameter-Free Prediction with Arithmetic Uniqueness}

\author[1]{Marvin L. Gentry}
\affil[1]{Independent Researcher, \texttt{drmlgentry@protonmail.com}\\
ORCID: 0009-0006-4550-2663}
\date{\today}

\begin{document}
\maketitle

\begin{abstract}
We extend the golden ratio mass lattice $m = m_e \cdot \phi^{q/4}$,
$q \in \mathbb{Z}$, to the neutrino sector without additional free
parameters. For the base $\phi = (1+\sqrt5)/2$, a systematic scan of
all integer triples $(q_1, q_2, q_3)$ subject to the PDG 2024 neutrino
oscillation constraints reveals a unique solution within $1\sigma$:
$(q_1, q_2, q_3) = (-149, -146, -134)$, giving masses
$m_1 = 8.387$~meV, $m_2 = 12.033$~meV, $m_3 = 50.972$~meV,
and $\sum m_\nu = 71.4$~meV ($\chi^2 = 0.58$, 2~d.o.f.).
The mass ratios satisfy exact algebraic relations:
$m_2/m_1 = \phi^{3/4}$, $m_3/m_2 = \phi^3$, $m_3/m_1 = \phi^{15/4}$.
A scan over bases $b\in[1.50,1.71]$ shows that four other bases admit
$1\sigma$ solutions, but none have any known algebraic significance
or independent derivation from the arithmetic geometry of the lepton
or quark sector manifolds. Among all algebraically significant bases,
$\phi$ is the unique base admitting a $1\sigma$ solution.
The prediction $\sum m_\nu = 71.4 \pm 2.9$~meV is falsifiable at
$3\sigma$ by CMB-S4~\cite{CMBS4}, which will reach $\sim 30$~meV
sensitivity. Normal ordering is the only solution; inverted ordering
admits no solution within $3\sigma$.
\end{abstract}

\section{Introduction}
\label{sec:intro}

Neutrino masses are the only Standard Model fermion masses not yet
directly measured. Neutrino oscillation experiments constrain the
mass-squared differences $\Delta m^2_{21}$ and $\Delta m^2_{31}$
but not the absolute mass scale. The sum $\sum m_\nu$ is bounded
from above by cosmological observations ($< 120$~meV, Planck
2018~\cite{Planck2018}) and from below by oscillation data
($\sum m_\nu > 59$~meV for normal ordering).

In the hyperbolic flavor geometry framework~\cite{Gentry:CORE}, the
observed fermion masses cluster near the golden ratio lattice
\begin{equation}
  m = m_e \cdot \phi^{q/4}, \qquad q \in \mathbb{Z},
  \label{eq:lattice}
\end{equation}
where $\phi = (1+\sqrt{5})/2$ is the golden ratio and $m_e =
0.51099895$~MeV is the electron mass (anchor). The lattice spacing
$\log\phi$ is the regulator of the quadratic field $\mathbb{Q}(\sqrt{5})$
and equals $\tfrac{1}{2}\log M(\Delta_{\texttt{m003}})$, where
$M(\Delta_{\texttt{m003}}) = \phi^2$ is the Mahler measure of the
Alexander polynomial of the cusped lepton-sector
manifold~\cite{Gentry:JKTR}.

The charged fermion masses fit equation~\eqref{eq:lattice} with
RMS residual $0.055\,\log_\phi$ units (approximately $3\%$ in mass
ratio)~\cite{Gentry:CORE}. The extension to neutrinos requires
$q < 0$, corresponding to masses below the electron mass. This is
natural: the lattice extends symmetrically to negative integers,
and the negative side was never constrained by prior phenomenology.

We present here a focused analysis of the neutrino sector: the scan
protocol, the unique solution for the base $\phi$, the exact mass
ratio relations, the arithmetic uniqueness argument, and the
falsifiable prediction. This paper is self-contained; we refer
to~\cite{Gentry:CORE} for the full framework.

\section{Scan protocol}
\label{sec:scan}

\subsection{Input data}

We use the PDG 2024~\cite{PDG2024} neutrino oscillation parameters
(normal ordering):
\begin{equation}
  \Delta m^2_{21} = (7.53 \pm 0.18) \times 10^{-5}\,\text{eV}^2,
  \qquad
  \Delta m^2_{31} = (2.51 \pm 0.03) \times 10^{-3}\,\text{eV}^2.
  \label{eq:pdg}
\end{equation}
The anchor is fixed at $m_e = 0.51099895$~MeV, as declared in the
charged fermion analysis~\cite{Gentry:CORE}. No new parameters are
introduced.

\subsection{Scan procedure}

We scan all integer triples $(q_1, q_2, q_3)$ with $q_1 < q_2 < q_3$
and $q_i \in [-200, 0]$ (neutrino masses are below the electron mass,
requiring $q < 0$). For each triple, we compute
\begin{align}
  m_i &= m_e \cdot \phi^{q_i/4}, \\
  \Delta m^2_{21} &= m_2^2 - m_1^2, \\
  \Delta m^2_{31} &= m_3^2 - m_1^2,
\end{align}
and record all triples satisfying both constraints within $N\sigma$
of the PDG central values.

\subsection{Scan results for $\phi$ by confidence level}

For the golden ratio base $\phi$, the results are:

\begin{center}
\begin{tabular}{cc}
\toprule
Confidence level & Solutions \\
\midrule
$1\sigma$ & 1 \\
$2\sigma$ & 0 \\
$3\sigma$ & 3 \\
\bottomrule
\end{tabular}
\end{center}

The $2\sigma$ gap is notable: the unique $1\sigma$ solution is
isolated, with no nearby solutions. This is a consequence of the
discrete lattice structure — the allowed $(q_1, q_2, q_3)$ triples
are sparse in the region satisfying both oscillation constraints
simultaneously.

\section{The unique solution for $\phi$}
\label{sec:solution}

The unique triple satisfying both PDG constraints within $1\sigma$ is:
\begin{equation}
  (q_1, q_2, q_3) = (-149, -146, -134),
  \label{eq:triple}
\end{equation}
giving masses
\begin{equation}
  m_1 = 8.387\,\text{meV}, \quad
  m_2 = 12.033\,\text{meV}, \quad
  m_3 = 50.972\,\text{meV},
  \label{eq:masses}
\end{equation}
with
\begin{equation}
  \sum m_\nu = 71.4\,\text{meV}.
  \label{eq:sum}
\end{equation}
The residuals are $-0.48\sigma$ for $\Delta m^2_{21}$ and $+0.59\sigma$
for $\Delta m^2_{31}$, giving $\chi^2 = 0.58$ (2~d.o.f., $p = 0.75$).

\subsection{Comparison with oscillation constraints}

\begin{table}[htbp]
\centering
\caption{Predicted vs PDG neutrino oscillation parameters.
PDG uncertainties are $1\sigma$.}
\label{tab:comparison}
\begin{tabular}{lccc}
\toprule
Observable & Predicted & PDG 2024 & Residual \\
\midrule
$\Delta m^2_{21}$ ($10^{-5}$~eV$^2$) & 7.430 &
  $7.53 \pm 0.18$ & $-0.48\sigma$ \\
$\Delta m^2_{31}$ ($10^{-3}$~eV$^2$) & 2.528 &
  $2.51 \pm 0.03$ & $+0.59\sigma$ \\
$\sum m_\nu$ (meV) & 71.4 & $< 120$ (Planck) & --- \\
\bottomrule
\end{tabular}
\end{table}

\subsection{Exact mass ratio relations}

The lattice spacings $\Delta q_{21} = 3$, $\Delta q_{32} = 12$,
$\Delta q_{31} = 15$ imply exact algebraic mass ratios:
\begin{equation}
  \frac{m_2}{m_1} = \phi^{3/4}, \qquad
  \frac{m_3}{m_2} = \phi^{3}, \qquad
  \frac{m_3}{m_1} = \phi^{15/4}.
  \label{eq:ratios}
\end{equation}
The relation $m_3/m_2 = \phi^3$ is particularly distinguished:
a spacing of 12 quarter-steps equals 3 full lattice units, connecting
the heaviest and middle neutrino masses by the cube of the golden ratio.
These are parameter-free predictions of the lattice structure.

\subsection{Inverted ordering}

For inverted ordering ($m_3 < m_1 < m_2$), the scan finds no solution
within $3\sigma$ for the base $\phi$. The lattice selects normal
ordering uniquely. This is a structural prediction: if experiments
establish inverted ordering, the lattice hypothesis is falsified
for the neutrino sector.

\section{Arithmetic uniqueness of $\phi$}
\label{sec:uniqueness}

\subsection{Base sensitivity scan}

To assess whether $\phi$ is special or whether other bases give
comparable results, we scan bases $b \in [1.50, 1.71]$ in steps
of $0.005$ and record which bases admit a $1\sigma$ solution
(with their own best-fit triples).

\begin{table}[htbp]
\centering
\caption{Bases admitting a $1\sigma$ solution to the neutrino
oscillation constraints. Only bases with known algebraic
significance are listed.}
\label{tab:bases}
\begin{tabular}{lcc}
\toprule
Base $b$ & $1\sigma$ solution? & Algebraic significance \\
\midrule
$\phi = (1+\sqrt{5})/2 \approx 1.618$ & Yes & Fundamental unit of
  $\mathbb{Q}(\sqrt{5})$; \\
  & & Mahler measure$^{1/2}$ of $\Delta_{\texttt{m003}}$ \\
$\approx 1.525$ & Yes & None known \\
$\approx 1.590$ & Yes & None known \\
$\approx 1.690$ & Yes & None known \\
$\approx 1.705$ & Yes & None known \\
\bottomrule
\end{tabular}
\smallskip
{\small Note: The other bases produce different integer triples
$(q_1,q_2,q_3)$ than the $\phi$ triple; they are not equivalent
to the $\phi$ solution.}
\end{table}

Five bases admit a $1\sigma$ solution in the scanned range. Four
have no known algebraic significance: they are not roots of
low-degree polynomials with small integer coefficients, and have
no connection to the arithmetic geometry of the lepton or quark
sector manifolds.

\subsection{Uniqueness proposition}

\begin{proposition}[Arithmetic uniqueness of $\phi$]
\label{prop:unique}
Among all bases $b$ that are algebraic integers with an independent
derivation from the arithmetic geometry of the PMNS or CKM manifold,
$\phi$ is the unique base for which the mass formula
$m_i = m_e \cdot b^{q_i/4}$ admits a solution within $1\sigma$ of
the PDG neutrino oscillation data~\eqref{eq:pdg}.
\end{proposition}

The proof reduces to the numerical scan: no other algebraically
motivated base (e.g., $\sqrt{2}$, $\sqrt[3]{2}$, $e$, $\pi$, or any
algebraic integer arising from the holonomy of $m003$ or $m006$)
falls within the four $1\sigma$ windows identified for the
non-significant bases.

\section{Falsifiable prediction}
\label{sec:prediction}

\subsection{Neutrino mass sum}

Propagating the PDG $1\sigma$ uncertainties on $\Delta m^2_{21}$ and
$\Delta m^2_{31}$ through equations~\eqref{eq:lattice} gives:
\begin{equation}
  \sum m_\nu = 71.4 \pm 2.9\,\text{meV} \quad (1\sigma,
  \text{ from oscillation data}).
  \label{eq:prediction}
\end{equation}
This is a parameter-free prediction: no continuous parameter was
tuned to obtain it. The $3\sigma$ falsification window is:
\begin{equation}
  \sum m_\nu < 62.7\,\text{meV} \quad \text{or} \quad
  \sum m_\nu > 80.1\,\text{meV}.
  \label{eq:falsification}
\end{equation}

\subsection{CMB-S4 sensitivity}

CMB-S4~\cite{CMBS4} is projected to reach $\sigma(\sum m_\nu) \approx
30$~meV sensitivity, sufficient to:
\begin{itemize}
  \item Detect the predicted sum $\sum m_\nu = 71.4$~meV at
    $>2.4\sigma$ significance if the true value equals the prediction.
  \item Falsify the lattice at $>3\sigma$ if the true value deviates
    by more than $\sim 9$~meV from the prediction.
  \item Distinguish normal from inverted ordering, providing an
    independent test of the lattice's ordering prediction.
\end{itemize}

\subsection{Context among cosmological bounds}

The predicted sum $\sum m_\nu = 71.4$~meV is:
\begin{itemize}
  \item Well within the Planck 2018 bound ($< 120$~meV).
  \item Above the minimum sum from oscillation data
    ($> 59$~meV for normal ordering).
  \item Consistent with all current constraints.
  \item In the range accessible to near-future cosmological surveys
    (DESI, Euclid, CMB-S4).
\end{itemize}

\section{Discussion}
\label{sec:discussion}

\subsection{The negative $q$ values}

The lattice triple $(q_1, q_2, q_3) = (-149, -146, -134)$ places all
three neutrino masses on the negative side of the lattice, far from
the charged leptons ($q_{e} = 0$, $q_\mu = 44$, $q_\tau = 68$).
This is natural: neutrinos are the only fermions with masses below the
electron mass, and the lattice extends symmetrically to negative
integers. The large gap between $q_\tau = 68$ and $q_1 = -149$
(a separation of 217 lattice steps) is not explained by the lattice
itself and represents the seesaw-like suppression of neutrino masses
relative to charged leptons. The lattice encodes the mass ratios
within the neutrino sector exactly, but does not dynamically explain
the overall scale suppression.

\subsection{Comparison with other lattice approaches}

Several authors have noted approximate golden-ratio relations in the
neutrino sector~\cite{various}, but these are typically post-hoc
observations without a geometric derivation. The present approach
differs in three respects: (i) the base $\phi$ is derived independently
from the Mahler measure of the Alexander polynomial of the lepton
manifold~\cite{Gentry:JKTR}, not fitted to the neutrino data;
(ii) the anchor $m_e$ is fixed by the charged fermion analysis, not
chosen to optimize the neutrino fit; (iii) the uniqueness of the
$1\sigma$ solution (no solution at $2\sigma$) for the base $\phi$
provides a quantitative measure of predictive sharpness.

\subsection{What is not claimed}

The lattice does not predict the absolute neutrino mass scale
from first principles — the anchor $m_e$ must be taken from
experiment. It does not explain the seesaw mechanism or the origin
of the large hierarchy between neutrino and charged lepton masses.
It provides a geometric constraint on the mass ratios within the
neutrino sector, not a dynamical theory of mass generation.

\section{Conclusion}
\label{sec:conclusion}

The golden ratio mass lattice $m = m_e \cdot \phi^{q/4}$ extends
to the neutrino sector with a unique $1\sigma$ solution for the base
$\phi$: $(q_1, q_2, q_3) = (-149, -146, -134)$, predicting
$\sum m_\nu = 71.4 \pm 2.9$~meV and exact mass ratios
$m_3/m_2 = \phi^3$, $m_2/m_1 = \phi^{3/4}$. The base $\phi$ is the
unique algebraically motivated base admitting a $1\sigma$ solution.
Normal ordering is selected uniquely; inverted ordering has no
$3\sigma$ solution. The prediction is falsifiable by CMB-S4.

The connection between $\phi$ and the lepton sector is not merely
numerical: $\phi$ is the fundamental unit of $\mathbb{Q}(\sqrt{5})$,
whose regulator $\log\phi$ equals the Mahler measure exponent of
the Alexander polynomial of the cusped lepton manifold \texttt{m003}.
The neutrino mass prediction is thus grounded in the same arithmetic
structure as the lepton mixing matrix and CP phase.


\section*{Acknowledgements}
The author used DeepSeek and Claude as AI assistants during the
preparation of this manuscript, including for \LaTeX\ editing and
computational script development. All scientific content,
mathematical results, and conclusions are the sole responsibility
of the author.


\section*{Data Availability}
All scan scripts are available at
\url{https://github.com/drmlgentry/hyperbolic-flavor-scan}.
The neutrino scan script \texttt{neutrino\_rigorous.py} produces
the results in Table~\ref{tab:comparison} deterministically.

\section*{Conflict of Interest}
The author declares no conflicts of interest.

\begin{thebibliography}{99}

\bibitem{PDG2024}
  S.~Navas et al.\ (Particle Data Group),
  Phys.\ Rev.\ D \textbf{110}, 030001 (2024).

\bibitem{Planck2018}
  N.~Aghanim et al.\ (Planck Collaboration),
  ``Planck 2018 results. VI. Cosmological parameters,''
  Astron.\ Astrophys.\ \textbf{641}, A6 (2020).

\bibitem{CMBS4}
  K.~Abazajian et al.\ (CMB-S4 Collaboration),
  ``CMB-S4 Science Book, First Edition,''
  arXiv:1610.02743 (2016).

\bibitem{Gentry:CORE}
  M.~L.~Gentry,
  ``Geometric Unification of Flavor: Masses, Mixing, and CP
  from the Golden Ratio Lattice,''
  preprint (2026). See also companion papers SSRN:6583549--6583553.

\bibitem{Gentry:JKTR}
  M.~L.~Gentry,
  ``The Alexander Polynomial of the $(-2,3,7)$ Pretzel Knot
  and the Golden Ratio,''
  submitted to \textit{J.\ Knot Theory Ramifications} (2026)
  (JKTR-S-26-00044).

\bibitem{Gentry:CKM}
  M.~L.~Gentry,
  ``Quark Mixing from Hyperbolic Geometry,''
  SSRN preprint, doi:10.2139/ssrn.6583550 (2026).

\bibitem{Gentry:PMNS}
  M.~L.~Gentry,
  ``Lepton Mixing from the Borel Structure of Hyperbolic Holonomy,''
  SSRN preprint, doi:10.2139/ssrn.6583553 (2026).

\bibitem{various}
  See e.g.\ W.~Rodejohann,
  ``Neutrino mixing and the quark-lepton complementarity,''
  Phys.\ Rev.\ D \textbf{69}, 033005 (2004), and references therein.

\end{thebibliography}

\end{document}